\section{先问人，再定义任务}

如果研究对象与社区生活、教育和语言延续直接相关，模型团队不能仅凭 benchmark 想象需求。NLLB 先用访谈确认问题，再让母语者参与翻译与评测；这是课堂里最容易被模型细节遮蔽、却最重要的方法论转折。

\subsection{怎样知道自己解决的是真问题}

本节从覆盖目标回到问题定义。在展示任何数据管线之前，Fan 先提出研究判断：最重要的是确认团队正在解决真实问题，尤其当技术与人非常接近时。这个问题不是一句伦理口号，而是决定后续数据标准、成功条件和错误代价的第一步。

\lecturefigure{slide-06-real-problem.jpg}{在建模前确认问题对真实使用者成立}{Stanford 课堂视频 00:05:06}

\begin{knowledgebox}{读图：这是一道研究流程门，而不是开场修辞}
如果使用者主要担心教育中的语言流失，系统就不能只优化新闻翻译；如果社区使用多个脚本或本地变体，统一到单一标准可能直接破坏可用性。需求访谈因此会改变数据采样、标签体系、评测场景和 release policy。
\end{knowledgebox}

\begin{importantbox}{Human-centered 不等于最后做一次用户测试}
更完整的参与链路是：
\[
\text{需求访谈}\rightarrow\text{语言标准协商}\rightarrow\text{专业翻译}
\rightarrow\text{母语评测}\rightarrow\text{错误反馈与修正}.
\]
使用者不是流水线末端的验收员，而是问题定义和证据生产的共同参与者。
\end{importantbox}

\subsection{访谈样本与方法边界}

团队访谈了 44 位母语者，涉及 36 种语言与多个地区。这个样本不能代表所有低资源语言社区，却足以暴露纯技术路线看不到的问题：语言衰退感、主流教育语言挤压、被现有翻译系统纳入的期待，以及“并不完美但足够有用”的现实需求。

\lecturefigure{slide-07-speaker-interviews.jpg}{低资源语言母语者访谈的范围}{Stanford 课堂视频 00:05:54}

\begin{knowledgebox}{读图：样本数字说明探索范围，不提供总体推断}
44 位受访者与 36 种语言显示团队主动跨语言收集意见，但每种语言对应的人数很少，且课堂也承认招募渠道带有地域与移民群体偏差。正确用途是形成 design hypotheses，再在翻译、评测和发布环节持续咨询，而不是把访谈比例写成全球统计结论。
\end{knowledgebox}

\teachervoice{Fan 把这一步称为 social-sciences-type 或 sociology-type approach。她不是说访谈能替代实验，而是说在贴近人的问题上，先把目标问清楚，才知道后续实验应该测什么。}

\subsection{语言衰退、技术纳入与“足够好”}

访谈结果把机器翻译放进了更大的社会背景：部分使用者担心本地语言在教育和公共生活中衰退；与此同时，被主流翻译系统纳入会带来强烈期待。课堂还指出，使用者并不一定要求机器翻译达到专业人工翻译的全部标准，某些日常沟通场景中，明显有缺陷但方向正确的工具也可能产生价值。

\lecturefigure{slide-08-interview-findings.jpg}{访谈揭示的担忧、期待与可用性判断}{Stanford 课堂视频 00:07:03}

\begin{knowledgebox}{读图：三条发现分别作用于三个系统层面}
语言衰退影响项目优先级与文化风险；“被纳入系统”的期待影响覆盖目标；自动翻译在现实使用中可能已足够，则要求评测区分专业出版、教育、旅行、健康等场景。不能用一个平均分替代这些不同风险容忍度。
\end{knowledgebox}

\begin{warningbox}{“足够有用”不是降低低资源语言标准}
场景化门槛允许不同任务有不同容错率，但不能据此把低资源语言永远停留在较差质量。尤其在医疗、公共卫生、法律和危机信息中，小词误译也可能造成严重后果，后文 toxicity 例子会具体展示这一点。
\end{warningbox}

\teachervoice{课堂 Q\&A 再次强调，母语者咨询应贯穿整个开发流程：从初始访谈，到专业翻译团队，再到人工评测与开放社区反馈。Human-centered 是生命周期属性，不是一次性 workshop。}

\subsection{本章小结}

NLLB 在建模前先确认真实需求，并把社区参与扩展到数据与评测。访谈提供的是设计证据和风险线索，不是代表性统计；“有用”必须按场景定义，不能成为降低质量责任的借口。

